%%%%%%%%%%%%%%%%%%%%%%%%%%%%%%%%%%%%%%%%%%%%%%%%%%%%%%%%
%%%% Las Variables siempre deben ir inclinadas      %%%%
%%%% los símbolos o subíndices van erguidas         %%%%
%%%%%%%%%%%%%%%%%%%%%%%%%%%%%%%%%%%%%%%%%%%%%%%%%%%%%%%%

%%%%%%%%%%%%%%%%%%%%%%  Masa - inicio %%%%%%%%%%%%%%%%%%%%%%%%%%
% 1. La base: #1=subíndice, #2=punto u otro acento (opcional)
% Significado: \comando[opcional]{obligatorio}
% Explicación:
% [2]: Indica que el comando recibe dos argumentos en total.
% [] (el segundo par de corchetes): Define que el primer argumento 
% es opcional y que, por defecto, está vacío.
% En LaTeX, el argumento opcional siempre es el #1 y se pasa entre 
% corchetes [].
% El argumento obligatorio es el #2 y se pasa entre llaves {}.
% \mass*[acento]{subindice}[exponente]
% Default: m  |  With *: M
\NewDocumentCommand{\mass}{s O{} m O{}}{%
  \ensuremath{%
    \let\baseM\relax
    % If * is used, switch to uppercase M; otherwise, use lowercase m
    \IfBooleanTF{#1}{\def\baseM{M}}{\def\baseM{m}}%
    \ifx\\#2\\ % If no accent is provided
      {\baseM}_{\text{#3}}^{#4}%
    \else % If an accent (\dot, \bar, etc.) is provided
      {#2{\baseM}}_{\text{#3}}^{#4}%
    \fi
  }%
}
\newcommand{\vaporMass}{\mass{\vapor}}
\newcommand{\gasMass}{\mass{\gas}}
\newcommand{\vaporFlowMass}[1][\dot]{\mass[#1]{\vapor}}   %%%% Flujo de masa del vapor: kg/s
\newcommand{\gasFlowMass}[1][\dot]{\mass[#1]{\gas}}       %%%% Flujo de masa del gas: kg/s
%%%%%%%%%%%%%%%%%%%%%%  fin - Masa  %%%%%%%%%%%%%%%%%%%%%%%%%%%

\newcommand{\molecularWeight}[2][]{%
  \ensuremath{{M}_{\text{#2}}^{#1}}%
}

\newcommand{\vaporMolecularWeight}{\molecularWeight[]{\vapor}}        %%%% Molecular weight vapor: M_vap
\newcommand{\gasMolecularWeight}{\molecularWeight[]{\gas}}            %%%% Molecular weight gas: M_g
\newcommand{\averageMolecularWeight}{\molecularWeight[]{}}            %%%% Molecular weight average

%%%%%%%%%%%%%%%%%%%%%%%%%%  MOLES para el estado gaseoso- inicio %%%%%%%%%%%%%%%%%%%%%%%%%%%%%%
\newcommand{\mole}[2][]{%
  \ensuremath{{n}_{\text{#2}}^{#1}}%
}
\newcommand{\vaporMole}{\mole[]{\vapor}}      %%%% Moles del vapor
\newcommand{\gasMole}{\mole[]{\gas}}          %%%% Moles del gas
\newcommand{\totalMole}{\mole[]{\total}}      %%%% Moles totales de la mezcla


\newcommand{\moleFractionPhaseGaseous}[2][]{%
  \ensuremath{{y}_{\text{#2}}^{#1}}%
}
\newcommand{\vaporMoleFraction}{\moleFractionPhaseGaseous[]{\vapor}}  %%%% Moles del vapor
\newcommand{\gasMoleFraction}{\moleFractionPhaseGaseous[]{\gas}}      %%%% Moles del gas
%%%%%%%%%%%%%%%%%%%%%%  fin - MOLES  %%%%%%%%%%%%%%%%%%%%%%%%%%%

%%%%%%%%%%%%%%%%%%%%%%%%%%  Presión - inicio %%%%%%%%%%%%%%%%%%%%%%%%%%%%%%
%%%% [2]: Indica que el comando acepta dos argumentos.
%%%% []: Indica que el primer argumento (#1) es opcional y que, por defecto, está vacío.
%%%% #1 is Optional: Because you put [] after the [2]. If you don't provide it, it defaults to "empty".
%%%% #2 is Mandatory: You must provide the {subscript} or LaTeX will grab the next character in your text as the argument.
%%%% \pressure{} produces: p
%%%% \pressure{v} produces: p_v
%%%% \pressure[]{} produces: p
%%%% \pressure[o]{v} produces: p_{v}^{o}
\newcommand{\pressure}[2][]{%
  \ensuremath{p_{\text{#2}}^{#1}}%
}
\newcommand{\pAtm}{\pressure{\atmospheric}}               % Atmospheric pressure
\newcommand{\pSat}{\pressure[\saturated]{\vapor}}         % Saturated pressure
\newcommand{\pSatInv}{{\{\pressure[{\saturated}]{\vapor}\}}^{-1}}  
\newcommand{\pVapor}{\pressure[]{\vapor}}                 % Vapor pressure (reuses your \vapor command)
\newcommand{\pGas}{\pressure{\gas}}                       % Gas pressure (reuses your \gas command)
\newcommand{\pTotal}{\pressure{\total}}                   % Presión total
\newcommand{\pCritic}{\pressure{\critical}}               % Presión crítica
\newcommand{\pTriple}{\pressure{\triple}}                 % Presión en el punto triple
% \newcommand{\pSat}{\pressure[o]{v}\xspace}
% \newcommand{\pVapor}{\pressure{\vapor}\xspace}
%%%%%%%%%%%%%%%%%%%%%%%%%%%  fin - Presión %%%%%%%%%%%%%%%%%%%%%%%%%%%%%%%


%%%% Temperatura
%%%% [2]: Indica que el comando acepta dos argumentos.
%%%% []: Indica que el primer argumento (#1) es opcional y que, por defecto, está vacío.
\newcommand{\temperature}[2][]{%
  \ensuremath{{T}_{\text{#2}}^{#1}}%
}
\newcommand{\tVapor}{\temperature{\vapor}}                    % Temperatura del Vapor
\newcommand{\tGas}{\temperature{\gas}}                    % Temperatura del Gas
\newcommand{\tAtm}{\temperature{\atmospheric}}   
\newcommand{\tCritic}{\temperature{\critical}}                    % Temperatura crítica
\newcommand{\tTriple}{\temperature{\triple}}                    % Temperatura en el punto triple
\newcommand{\tDewPoint}{\temperature{\dewPointSymbol}}                    % Temperatura en el punto rocío

%%%% Generic Volume
% \volume*[acento]{subindice}[exponente]
\NewDocumentCommand{\volume}{s O{} m O{}}{%
  \ensuremath{%
    \let\baseV\relax
    \IfBooleanTF{#1}{\def\baseV{v}}{\def\baseV{V}}%
    \ifx\\#2\\ % Si no hay acento
      {\baseV}_{\text{#3}}^{#4}%
    \else % Si hay acento (\dot, \bar, etc)
      {#2{\baseV}}_{\text{#3}}^{#4}%
    \fi
  }%
}
\newcommand{\specificVolume}{\volume*{}}                %%%% Volumen específico (minúscula v)
\newcommand{\vaporVolume}{\volume{v}}                   %%%% Volumen de vapor (V_v)
\newcommand{\gasVolume}{\volume{g}}                     %%%% Volumen de gas (V_g)
\newcommand{\fluxVolume}{\volume[\dot]{}}               %%%% Flujo  volumétrico (V con punto)
\newcommand{\molarVolume}{\volume[\bar]{}}            %%%% Volumen molar (V con barra)


%%%%%%%%%%%%%%%%%%%%%%  Heat - inicio %%%%%%%%%%%%%%%%%%%%%%%%%%
% 1. La base: #1=subíndice, #2=punto u otro acento (opcional)
% Significado: \comando[opcional]{obligatorio}
% Explicación:
% [2]: Indica que el comando recibe dos argumentos en total.
% [] (el segundo par de corchetes): Define que el primer argumento 
% es opcional y que, por defecto, está vacío.
% En LaTeX, el argumento opcional siempre es el #1 y se pasa entre 
% corchetes [].
% El argumento obligatorio es el #2 y se pasa entre llaves {}.
% \heat*[acento]{subindice}[exponente]
% Default: Q  |  With *: q
\NewDocumentCommand{\heat}{s O{} m O{}}{%
  \ensuremath{%
    \let\baseQ\relax
    % If * is used, switch to uppercase M; otherwise, use lowercase m
    \IfBooleanTF{#1}{\def\baseQ{q}}{\def\baseQ{Q}}%
    \ifx\\#2\\ % If no accent is provided
      {\baseQ}_{\text{#3}}^{#4}%
    \else % If an accent (\dot, \bar, etc.) is provided
      {#2{\baseQ}}_{\text{#3}}^{#4}%
    \fi
  }%
}
\newcommand{\heatFlow}{\heat[\dot]{}}  
%%%%%%%%%%%%%%%%%%%%%%%%%%%  fin - Heat %%%%%%%%%%%%%%%%%%%%%%%%%%%%%%%

%%%%%%%%%%%%%%%%%%%%%%  Specific Heat - inicio %%%%%%%%%%%%%%%%%%%%%%%%%%
% 1. La base: #1=subíndice, #2=punto u otro acento (opcional)
% Significado: \comando[opcional]{obligatorio}
% Explicación:
% [2]: Indica que el comando recibe dos argumentos en total.
% [] (el segundo par de corchetes): Define que el primer argumento 
% es opcional y que, por defecto, está vacío.
% En LaTeX, el argumento opcional siempre es el #1 y se pasa entre 
% corchetes [].
% El argumento obligatorio es el #2 y se pasa entre llaves {}.
% \specificHeat*[acento]{subindice}[exponente]
% Default: C  |  With *: c
\NewDocumentCommand{\specificHeat}{s O{} m O{}}{%
  \ensuremath{%
    \let\baseC\relax
    \IfBooleanTF{#1}{\def\baseC{c}}{\def\baseC{C}}%
    \ifx\\#2\\ % If no accent is provided
      {\baseC}_{\text{#3}}^{#4}%
    \else % If an accent (\dot, \bar, etc.) is provided
      {#2{\baseC}}_{\text{#3}}^{#4}%
    \fi
  }%
}
\newcommand{\specificHeatAtV}{\specificHeat[]{v}}  
\newcommand{\specificHeatAtP}{\specificHeat[]{p}}  
%%%%%%%%%%%%%%%%%%%%%%%%%%%  fin - Specific Heat %%%%%%%%%%%%%%%%%%%%%%%%%%%%%%%




%%%% Representación para Humedad
\newcommand{\humidity}[2][]{%
  \ensuremath{{Y}_{\text{#2}}^{#1}}%
}
\newcommand{\absoluteHumidity}{\humidity[]{}}                     %%%% Humedad absoluta
\newcommand{\molarAbsoluteHumidity}{\humidity[]{\molar}}          %%%% Humedad absoluta molar
\newcommand{\saturatedAbsoluteHumidity}{\humidity[]{\saturated}}  %%%% Humedad absoluta saturada
\newcommand{\specificHumidity}{\humidity{w}}                      %%%% Humedad específica
\newcommand{\volumetricHumidity}{\humidity{w}}                    %%%% Humedad volumétrica

%%%% Variable para la humedad relativa
\newcommand{\relativeHumidity}[2][]{%
  \ensuremath{{\phi}_{\text{#2}}^{#1}}%
}

%%%% Factor acéntrico
\newcommand{\acentricFactor}[2][]{%
  \ensuremath{{\omega}_{\text{#2}}^{#1}}%
}

%%%% Entalpía: H
\newcommand{\enthalpy}[2][]{%
  \ensuremath{{H}_{\text{#2}}^{#1}}%
}
\newcommand{\specificEnthapy}{{{h}}}
% \newcommand{\diffEnthalpy}{\heat{\latent{}}}

%%%% Cambio de entalpía de vaporización
\newcommand{\dhVap}{\Delta \enthalpy[]{\vapor}}



\newcommand{\entropy}[2][]{%
  \ensuremath{{S}_{\text{#2}}^{#1}}%
}

\newcommand{\gibbsEnergy}[2][]{%
  \ensuremath{{G}_{\text{#2}}^{#1}}%
}

\newcommand{\helmotzEnergy}[2][]{%
  \ensuremath{{A}_{\text{#2}}^{#1}}%
}

%%%% Calor total                                        %%%%
\newcommand{\totalHeat}{\heat{\total{}}}

%%%% Calor sensible
\newcommand{\sensibleHeat}{\heat{\sensible{}}}

%%%% Calor latente
\newcommand{\latentHeat}{\heat{\latent{}}}

\newcommand{\heatCapacityAtConstantPressure}{\heatCapacity{p}}

\newcommand{\heatCapacityAtConstantVolume}{\heatCapacity{v}}






%%%%


%%%%


%%%%
\def\psychrometricRatio{\ensuremath{{r}}}

%%%% Coeficiente de saturación
\newcommand{\degreeOfSaturation}[2][]{%
  \ensuremath{{\mu}_{\text{#2}}^{#1}}%
}
% Factor de compresibilidad
\newcommand{\factorCompressibility}[2][]{%
  \ensuremath{{z}_{\text{#2}}^{#1}}%
}
\newcommand{\zVapor}{\factorCompressibility[]{\vapor}}  %%%% Factor de compresibilidad para el vapor
\newcommand{\zGas}{\factorCompressibility[]{\gas}}      %%%% Factor de compresibilidad para el gas
